\documentclass[11pt]{article}

\newcommand{\cnum}{Math 245B}
\newcommand{\ced}{Winter 2026}
\newcommand{\ctitle}[4]{\title{\vspace{-0.5in}\cnum, \ced\\Week #1: #2}\author{\vspace{-0.35in}\\#3, #4}}
\usepackage{enumitem}
\usepackage[usenames,dvipsnames,svgnames,table,hyperref]{xcolor}
\usepackage{amsmath, amsfonts}
\usepackage{graphicx} % Required for inserting images
\usepackage{float}
\usepackage{listings}
\usepackage{xcolor}
\usepackage{hyperref}
\usepackage{comment}

\renewcommand*{\theenumi}{\alph{enumi}}
\renewcommand*\labelenumi{(\theenumi)}
\renewcommand*{\theenumii}{\roman{enumii}}
\renewcommand*\labelenumii{\theenumii.}

\definecolor{codegreen}{rgb}{0,0.6,0}
\definecolor{codegray}{rgb}{0.5,0.5,0.5}
\definecolor{codepurple}{rgb}{0.58,0,0.82}
\definecolor{backcolour}{rgb}{0.95,0.95,0.92}
\definecolor{header}{HTML}{205459}
\definecolor{solution}{HTML}{204d52}
\DeclareMathOperator{\spt}{spt}
\DeclareMathOperator{\iter}{int}

\newcommand{\solution}[1]{{{\textcolor{header}{Solution:} \textcolor{solution}{#1}}}}
\setlist[enumerate]{label=(\alph*)}

\lstdefinestyle{mystyle}{
    backgroundcolor=\color{backcolour},
    commentstyle=\color{codegreen},
    keywordstyle=\color{magenta},
    numberstyle=\tiny\color{codegray},
    stringstyle=\color{codepurple},
    basicstyle=\ttfamily\footnotesize,
    breakatwhitespace=false,
    breaklines=true,
    captionpos=b,
    keepspaces=true,
    numbers=left,
    numbersep=5pt,
    showspaces=false,
    showstringspaces=false,
    showtabs=false,
    tabsize=2
}


\begin{document}
\ctitle{\#2}{}{Asher Christian}{006-150-286}
\date{}
\maketitle

\section{Monday}
Regularization of $L^{p}$-function.\\
Notation: $\Omega \subset \mathbb{R}^{d}$ and for $\delta > 0$ $\Omega_\delta = \{ x \in \Omega : d(x,\Omega^{c}) > \delta\}$
\begin{enumerate}
    \item Given $f : \Omega \rightarrow \mathbb{R}$, $h \in \mathbb{R}^{d}$, we define $\tau_h(F) : \Omega_|h| \rightarrow \mathbb{R}$ by
        \[
        \tau_h(f)(x) = f(x + h)
        \] 
    \item If $w \subset \mathbb{R}^{d}$, we write $\omega \subset \subset \Omega$ if $\omega$ is bounded
        and  $\overline{ \omega} \subset \Omega$ 
    \item If $K \subset \Omega$ is compact, we define the $L^{p}$-variations of $f$ over  $K$ by
         \[
             ||\tau_h(f) - f||_{L^{p}(K)} \qquad K \subset \Omega_{|h|}
        \] 
\end{enumerate}
Remark:  
\[
    ||\tau_h(f) - f||_{L^{p}(K)^{p}} = \int_{K}^{}|f(x+h) - f(x)|^{p}dx
\] 
therefore
\[
    \frac{||\tau_h(f) - f)||_{L^{p}}^{p}}{|h|^{p}} = \int_{K}^{} |\frac{f(x+h) - f(x)}{|h|}|^{p} dx
\] 
\[
    F \subset C(K) \text{ is precompact if and only if $F$ is bounded and equicontinuous}
\] 
Notation: We define $\overline{ \rho} : \mathbb{R}^{d} \rightarrow [0, +\infty)$ by
\[
\overline{ \rho }(x) = 
\begin{cases}
    0 & |x|_{l_2} \ge 1 \\
    e^{-\frac{1}{1-|x|_{l_2}^2}} & |x|_{l_2} < 1
\end{cases}
\] 
we have that $\overline{ \rho} \ge 0, \overline{ \rho}(x) = 0$ if and only if $|x|_{l_2} \ge 1$
and $\overline{ \rho} \in C^{\infty}( \mathbb{R}^{d})$ we have $|x|_{l_2}^2 = x_1^2 + \dots + x_d^2$ this is the standard modifier function.
We set, $\rho = \frac{\overline{ \rho}}{|| \overline{ \rho}||_{L^{1}( \mathbb{R}^{d})}}$ so that
\[
\int_{ \mathbb{R}^{d}}^{} \rho(x) dx = 1
\] 
For $\epsilon > 0$ we define
\[
    \rho_\epsilon(y) = \frac{1}{\epsilon^{d}} \rho(\frac{y}{\epsilon}) , \rho_\epsilon \in C^{\infty}( \mathbb{R}^{d}), \rho_\epsilon(y) > 0 \iff |y|_{l_2} < \epsilon
\] 
\[
\int_{ \mathbb{R}^{d}}^{} \rho_\epsilon(y)dy = \int_{ \mathbb{R}^{d}}^{} \rho(\frac{y}{\epsilon}) d(\frac{y}{\epsilon}) = 1
\] 
and $\rho_0 = \delta_0$ since..
Convolution of two function. Let $f,g : \mathbb{R}^{d} \rightarrow [0,+\infty)$ measurable wrt lebesgue.
then
\[
f * g(x) = \int_{ \mathbb{R}^{d}}^{}f(x-y)g(y)dy = g * f(x)
\] 
if $f \in L^{p}( \mathbb{R}^{d}), g \in L^{q}( \mathbb{R}^{d})$ then $f * g \in L^{r} ( \mathbb{R}^{d})$ and $||f * g||_{L^{r}} \le ||f||_{L^{p}} ||g||_{L^{q}}$ where
\[
1 + \frac{1}{r} = \frac{1}{p} + \frac{1}{q}
\] 
when you compute
\[
    \lim_{\epsilon\to 0 } \rho_\epsilon * f = f \quad f \in C( \mathbb{R}^{d}), \rho_\epsilon * f \rightarrow f \text{ everywhere}
\] 
if $f \in L^{p} (  \mathbb{R}^{d})$ then $||\rho_\epsilon * f - f||_{L^{p}} \rightarrow 0$ and $\rho_\epsilon * f(x) \rightarrow f(x)$ if $x$ is a lebesgue point for  $x$.
We next state a variant of the arscoli arzela theorem for  $L^p$ function.
\section{Theorem (M. Riesz - Frechet-Kolmogorov)}
let  $1 \le p < \infty$ and let $\Omega \subset \mathbb{R}^{d}$ be an open set.
Let $\omega \subset \subset \Omega$ be bounded and let $F$ be a bounded subset of  $L^{p}(\Omega)$ satisfying the following property
\begin{enumerate}
    \item For every $\epsilon > 0$ there exists $\delta > 0$ such that if $h \in  \mathbb{R}^{d}$,  $|h|_{l_2} \le \delta$ then $||\tau_h(f) -f||_{L^{p}(\omega)} \le \epsilon$
\end{enumerate}
Then $F_{\omega} = \{f|_{\omega} : f \in F\} $ is precompact in  $L^{p}(\omega)$ (for the strong $L^{p}$ topology).\\
Proof: since $\omega$ is bounded, it is not a loss of generality to assume that  $\Omega$ is bounded. Since  $L^{p}(\omega)$ is complete
it suffices to show that the closure of $F_\omega$ in  $L^{p}(\omega)$ is totally bounded. Given $\epsilon > 0$ we are to find
\[
f_1,\dots,f_n \in F_\omega
\] 
such that
\[
    F_\omega \subset \bigcup_{i =1}^{n}B_\epsilon^{L^{p}(\omega)}(f_i)
\] 
choose $\delta > 0$ such that $(a)$ holds.
\begin{enumerate}
     \item Approximation by smooth functions
\end{enumerate}
We claim that if $n > \frac{1}{\delta}$ then $||\rho_{\frac{1}{n}} * f - f||_{L^{p}(\omega)} < \epsilon$.
We have since $\rho$ integrates to 1
\[
    |\rho_{\frac{1}{n}} * f(x) - f(x)|^{p} = |\int_{ \mathbb{R}^{d}}^{} \rho_{\frac{1}{n}}(y)(f(x-y) - f(x)) dy|^{p} \le \int_{ \mathbb{R}^{d}}^{} \rho_{\frac{1}{n}}(y) |f(x-y) -f(x)|^{p} dy
\] 
\[
    = \int_{ \mathbb{R}^{d}}^{} \rho_{\frac{1}{n}}(y)|\tau_{-y}f(x) - f(x)|^{p}dy 
\] 
thus
\[
    \int_{\omega}^{}|\rho_{\frac{1}{n}} * f(x) - f(x)|^{p}dx \le \int_{ \omega}^{}\int_{ \mathbb{R}^{d}}^{}\rho_{\frac{1}{n}}(y)|\tau_{-y}f(x)-f(x)|^{p}dydx = \int_{ B_{\frac{1}{n}}(0)}^{}\rho_{\frac{1}{n}}(y)dy \int_{ \omega}^{}|\tau_{-y}f(x)-f(y)|^{p}dx
\] 
on this ball $|y| \le \frac{1}{n} \le \delta$\\
Since for $|y| \le \frac{1}{n}$ we have $|y| \le \delta$ we use $(a)$ to infer that
 \[
     ||\rho_{\frac{1}{n}} * f - f||^{p}_{L^{p}} \le \int_{ B_{\frac{1}{n}}}^{}\rho_{\frac{1}{n}}(y)\epsilon^{p}dy = \epsilon^{p}
\] 

\section{Wednesday}
Reminder
\begin{enumerate}
    \item $\rho \in C_c^{\infty}( \mathbb{R}^{d})$, $0 \le \rho_\epsilon(y) = \frac{1}{\epsilon} \rho(\frac{y}{\epsilon})$ and $\rho_\epsilon > 0$ iff $|y|_{l_2} < \epsilon$
        \[
        \int_{ \mathbb{R}^{d}}^{}\rho_\epsilon(y)dy = 1
        \] 
    \item $1 \le p < +\infty$, $\Omega \subset \mathbb{R}^{d}$ is open, $\omega \subset \subset \Omega$ is bounded, $F \subset L^{p}(\Omega) $ is bounded
    \item For all $\epsilon > 0$ there exists $\delta \in (0, d(\omega,\Omega^{c}) : h \in \mathbb{R}^{d}, |h|_{l_2} \le \delta \implies ||\tau_h(f) - f||_{L^{p}(\omega)} \le \epsilon$ for all $f \in F$
\end{enumerate}
Then $F_\omega$ is precompact in  $L^{p}(\omega)$.\\
Proof: Fix $\epsilon > 0$. We plan to prove that there exists $f_1,\dots,f_k \in F_\omega$ such that $F_\omega \subset \bigcup_{j=1}^{k}B^{L^{p}(\omega)}_\epsilon(f_j)$
\begin{enumerate}
    \item Done: choose $\delta$ such that $(c)$ holds and  $n > \frac{1}{\delta}$ then
        \[
            (1b) \qquad ||\rho_{\frac{1}{n}} * f - f||_{L^{p}(\omega)} < \epsilon \quad \forall f \in F_\omega
        \] 
    \item Let $D = \sup_{f \in F} ||f||_{L^{p}(\Omega)}$
    \item Set $H_n = \{\rho_{\frac{1}{n}} * f : f \in F_\omega\}$ we claim that $H_n$ is bounded in  $C(\overline{\omega})$ and equicontinuous.
        Indeed if $x \in \omega$
         \[
             |\rho_{\frac{1}{n}} * f(x)|^{p} =\left| \int_{ B_{\frac{1}{n}}(x)}^{}\rho_{\frac{1}{n}}(x-y)f(y)dy\right|^{p} \le \int_{B_{\frac{1}{n}}(x)}^{a} \rho_{\frac{1}{n}}(x-y)|f(y)|^{p}dy
        \] 
        \[
            \le c_n ||f||_{L^{p}}^{p} \quad C_n = \sup_{x,y} \rho_{\frac{1}{n}}(x-y)
        \] 
        use the fact that
        \[
        \nabla( \rho * f) = (\nabla \rho) * f
        \] 
        similarly
        \[
            |\nabla (\rho_{\frac{1}{n}} * f)(x)| = |(\nabla \rho_{\frac{1}{n}}) * f(x)| \le D^{p} c_n^{'} \quad c_n^{'} = \sup_{ \mathbb{R}^{d}} | \nabla \rho_{\frac{1}{n}}|
        \] 
        By $2$ and $3$,  $H_n$ is bounded in   $C(\overline{\omega}) $ and equicontinuous.
        Thus by the Ascoli-Arzela theorem $H_n$ is precompact in  $ C(\overline{\omega})$. Hence, there exist $f_1,\dots,f_k \in F_\omega$ such taht
        \[
        H_n \subset \bigcup_{j=1}^{k}B^{L^{p}(\omega)}_\epsilon( \rho_{\frac{1}{n}} * f_1)
        \] 
        since the balls in $L^{p}$ are equivalent to the balls in $C(\omega)$ since the measure of  $\omega$ is finite
        We use $1b$ to conlucde that
         \[
             F_\omega \subset \bigcup_{j=1}^{k}B_{2\epsilon}^{L^{p}(\omega)}(\rho_{\frac{1}{n}} * f_j)
        \] 
        you can change the center and increase the radius of the ball if these functions are not in $F_\omega$.
        Hence $\overline{ \mathcal{F}^{L^{p}(\omega)}}$ is totally bounded and complete hence compact.

\end{enumerate}

\section{Exercise}
Let $\Omega \subset \mathbb{R}^{d}$ be an open bounded set. Suppose that
\begin{enumerate}
    \item $\forall \epsilon > 0, \omega \subset \subset \Omega$ open there exists $\delta \in (0, d(\omega, \Omega^{c}))$ such that $h \in \mathbb{R}^{d}$, $|h| < \delta$
        implies
        \[
            ||\tau_h(f) - f||_{L^{p}(\omega)} \le \epsilon
        \] 
    \item for all $\epsilon > 0$ there exists $\omega \subset \subset \Omega$ open such that $||f||_{L^{p}(\Omega-w)} \le \epsilon$
    \item Show that $F$ is precompact in $L^{p}(\Omega)$
\end{enumerate}

\section{Hilbert spaces}
Definition: Let $H$ be a vector space. We call $\langle \cdot, \cdot \rangle : H^2 \rightarrow \mathbb{R}$ an inner product
if
\begin{enumerate}
    \item $\langle \lambda_1x_1 + \lambda_2 x_2, y \rangle = \lambda_1 \langle x_1,y \rangle + \lambda_2 \langle x_2, y \rangle$
    \item $\langle x_1,y \rangle = \langle y,x_1 \rangle$ 
    \item $\langle x, x \rangle >0$ for $x \in H \setminus \{0_H\}$
\end{enumerate}
for all $x_1,x_2,y \in H$.
In this case $x \to \sqrt{ \langle x, x \rangle} = ||x||$ is a norm .\\
Def: If $\langle \cdot, \cdot \rangle$ is an inner product on a vector space $H$ and $(H, ||\cdot||)$ is complete
we call  $(H, \langle \cdot, \cdot \rangle)$ a Hildber space.\\
Definition: A Banach space E is uniformly convex if for all $\delta > 0$ there exists $\epsilon > 0$ such that
if $x,y \in B_1^{E}(0)$ and $||x-y||_E \ge \delta$ then $|| \frac{x+y}{2}||_E  \le 1- \epsilon$.\\
Theorem: If $B$ is a uniformly convex Banach space then $B$ is reflexive.
Uniform convexity is a geometric property and reflexivity is a topological property

\section{Friday}
Theorem $H $ is a hilbert space implies  $H$ is reflexive. It suffices tos how that  $H$ is uniformly convex. Let  $x,y \in B_H$ then
 \[
||\frac{x+y}{2}||^2 + || \frac{x-y}{2}||^2 = \frac{||x||^2 + ||y||^2}{2} \le 1
\] 
if $\epsilon > 0$ and $||x-y|| \ge \epsilon$ then
\[
||\frac{x+y}{2}||^2 + \frac{\epsilon^2}{4} \le 1  \qquad || \frac{x+y}{2} ||^2 \le \sqrt{1 - \frac{\epsilon^2}{4}} = 1- \delta
\] 
for $0 < \delta < 1$.
Cauchy-Schwarz inequality: We have
\[
|\langle x, y \rangle | \le ||x|| ||y||
\] 
for all $x,y \in H$
 \[
0 \le ||x \pm y ||^2 = ||x||^2 + ||y||^2 + 2 \langle x, y \rangle
\] 
so
\[
| \langle x,y \rangle| \le \frac{||x||^2 + ||y||^2}{2}
\] 
\[
| \langle x ,y \rangle| = | \langle \lambda x, \frac{y}{\lambda} \rangle | \le \frac{|| \lambda x||^2 + \frac{1}{ \lambda} ||y||^2}{2} = ||x|| ||y||
\] 
Let $C \subset H$ be a non empty closed convex set. Then for every $f \in H$ there exists a unique  $ u \in C$ such that  $||f - u|| \le ||f - v|| $for all  $v \in C$\\
We set $u = P_C(f)$. Furthermore,  $u$ is characterized by
 \[
     (0) \quad
         \begin{cases}
             \langle f - u, v - u \rangle \le 0 & \forall v \in C\\
             u \in C
         \end{cases}
\] 
Proof: 1. Existence. Let $(x_n)_n \subset C$ be such that $||f - x_n|| \le D + \frac{1}{n}$ where $D = \inf_{v \in C} || f - v||$
Replacing $x_n -f$ by  $x_n$ we can assume that  $f = 0$. Since  $C$ is convex.
 \[
D \le ||\frac{x_n + x_m}{2}|| \le \frac{1}{2}( ||x_n|| + ||x_m||) \le D + \frac{1}{2n} + \frac{1}{2m}
\] 
so
\[
    (1) \quad \lim_{n,m \to \infty} || \frac{ X_n + x_m}{2} || = D
\] 
Recall that
\[
    (2) \quad || \frac{x_n + x_m}{2}||^2  + || \frac{x_n - x_m}{2} ||^2 = \frac{||x_n||^2 + ||x_m||^2}{2} \le \frac{1}{2}(D + \frac{1}{n})^2 + \frac{1}{2}(D + \frac{1}{m})^2
\] 
using $(1)$ and $(2)$ we infer
 \[
\lim_{n,m\to \infty} || \frac{x_n - x_m}{2} || = 0
\] 
and so $(x_n)$ is a cauchy sequence. Since  $H$ is complete $u := \lim_{n\to \infty}x_n$ exists.
Since $C$ is closed and $x_n \in C$ we have that  $u \in C$.
By assumption 
 \[
D = \lim_{n\to \infty}|| f - x_n|| = || f- u||
\] 
which proves the minimality property. We postpone the proof of uniqueness to a stronger corollary.
Let $v \in C$ and set  $N(t) = ||  (1-t)u + tv - f||^2 = || u + t(v-u) - f||^2$ for $t \in [0,1]$
Since  $(1-t)u + tv \in C$ we have  $N(t) \ge N(0) = D$
\[
||u-f||^2 + t^2 ||v -u||^2 + 2t \langle u -f, v- u \rangle
\] 
thus
\[
0 \le \lim_{t\to 0^{+}} \frac{N(t) - N(0)}{t} = 2 \langle u - f, v -u \rangle
\] 
hence
\[
\langle f -u , v - u \rangle \le 0
\] 
conversely, suppose that $(0)$ holds and let  $v \in C$. THen
 \[
||v- f||^2 = || v- u + u - f||^2 = || v- u||^2 + || u - f||^2 + 2 \langle v - u, u - f \rangle > ||u - f||^2
\] 
unless $ v = 0$.
Further assume  $C$ is a vector space.
Then 
\[
    (0) \iff \langle f -u, v \rangle \le \langle f - u,u \rangle \quad \forall v \in C
\] 
\[
    \implies \langle f - u, v \rangle = 0 \quad \forall v \in C
\] 
Theorem:
Suppose that $C \subset H$ is closed convex and nonempty. Then $P_c$ is 1-lipschitz.
 \[
||P_c(f_1) - P_c(f_2) || \le ||f_1 - f_2|| \quad \forall f_1,f_2 \in H
\] 
By $(0)$ 
 \[
\langle f_1 - P_c(f_1) , v - P_c(f_1)\rangle \le 0 \quad \langle f_2 - P_c(f_2) , v - P_c(f_2) \rangle \le 0 \quad \forall v \in C
\] 
\[
\implies \langle f_1 - P_c(f_1) , P_c(f_2) - P_c(f_1) \rangle \le 0 \quad \langle f_2 - P_c(f_2) , P_c(f_1) - P_c(f_2) \rangle \le 0
\] 
adding up we have 
\[
\langle f_1 - P_c(f_1) - f_2 + P_c(f_2) , P_c(f_2) - P_c(f_1) \rangle \le 0 \iff  || P_c(f_1) - P_c(f_2) || \le \langle f_2 - f_1, P_c(f_2) - P_c(f_1) \rangle 
\] 
\[
\le ||f_2 - f_1|| || P_c(f_2) - P_c(f_1) ||
\] 

\section{Riesz representation}
Given $\phi \in H'$ there exiss a unique $u \in H$ such that  $\phi(v) = \langle v, u \rangle$ for all $v \in H$.
Proof: Set  $C = Ker(\phi)$ to avoid trivialities we assume $\phi \ne 0$ identically and so $C \ne H$
By the above remark,  $ \langle f - P_c(f), v \rangle = 0$ for all $ c \in C$ so  $f = f - P_c(f) + P_c(f)$ where  $f - P_C(f) \in C^{\perp}$.
If we define 
\[
    C^{\perp} = \{v \in H : \langle v, c \rangle = 0, \forall c \in C\}
\] 
then we have proved that  $H = C \oplus C^{\perp}$ and note that
\[
Ker(\phi) = Ker(Id - P_c)
\] 
By assumption $C^{\perp} \ne \{0\}$ let $z \in C^{\perp}$ be such that $||z|| = 1$.
For any  $v \in H$
 \[
v \phi(z)  - z \phi(v) \in C
\] 
and so
\[
    (3) \quad 0 = \langle v \phi(z) - z \phi(v), z \rangle = \phi(z) \langle v, z \rangle - ||z||^2 \phi(v) 
\] 
\[
\phi(z) \langle v,z \rangle = \phi(v)
\] 
setting $u = \frac{z}{\phi(z)}$ we use $(3)$ to infer
\[
    \langle v, u \rangle = \phi(v)
\] 
Uniqueness is a consequence of the fact that if $\langle u_2 - u_1, v \rangle = 0$ for all $v \in H$ then  $u_1 = u_2$


\end{document}
